% =============================================================================
% 03_METHODS — Pipeline Design & Systematic Review Methodology
% =============================================================================
% Page target: ~2 pages
% Structure: Pipeline Design → Search Strategy → Selection → Analysis
% =============================================================================

\section{Methods}
\label{sec:methods}
\addcontentsline{toc}{section}{Methods}

This section has two components. First, it describes the automated
literature collection pipeline built to support this review — the technical
and conceptual architecture that enabled systematic coverage of a rapidly
evolving literature. Second, it describes the systematic review methodology
itself: the search strategy, selection criteria, quality assessment, and
analytical framework applied to the gathered corpus.

\subsection{Pipeline Design and Implementation}
\label{sec:pipeline}

The exponential growth of health AI literature between 2020 and 2026 rendered
the traditional one-time database search inadequate for a comprehensive
review. The volume of relevant publications — spanning computer science,
clinical medicine, law, ethics, and implementation science — exceeds what
any single research team can manually track. To address this, an automated
daily literature collection and triage pipeline was designed and deployed.

\subsubsection{Architecture Overview}

The pipeline operates across five layers: \textbf{data collection},
\textbf{paper scoring}, \textbf{LLM-based summarization},
\textbf{storage and indexing}, and \textbf{dissemination}.

Data collection is performed by \texttt{collector.py}, which queries multiple
bibliographic sources on a daily schedule: arXiv (preprints from cs.AI,
cs.LG, cs.CY, q-bio, eess.AS, and cs.CV), PubMed and PubMed Central (MEDLINE
via E-utilities), Semantic Scholar, and DBLP (capturing top-tier AI-ethics
venues including FAccT, AIES, and NeurIPS). Boolean search queries are
constructed dynamically around the five thematic keyword groups defined
for this review: AI Evaluation and Validation, Participatory Governance,
Adaptive Regulation, Evidence and Implementation, and Clinical Health AI.
Each query combination (e.g., \texttt{"AI evaluation" AND "healthcare"},
\texttt{"participatory governance" AND "clinical AI"}) is submitted to each
source, and results are aggregated with deduplication by DOI.

Each collected paper is scored by \texttt{ranker.py} across six dimensions:
theme relevance (composite score across the five thematic keyword groups),
novelty (publication recency), venue prestige, methodological rigour, and
an overall composite score. Papers scoring above a priority threshold are
designated Priority 3 (highest), with the remainder split across Priority
2 and Priority 1. This scoring enables focus of human review attention on
the highest-value publications while maintaining systematic coverage of the
literature landscape.

Priority 3 papers are passed to \texttt{summarizer.py}, which uses a
large language model (Gemini) to generate a structured summary for each
paper. Each summary encodes: title, authors, year, source, country, AI
system and method, evaluation method, governance dimension, main finding,
thematic tags, priority score, and a free-text justification of relevance.
These summaries are stored in both a JSON document store
(\texttt{data/papers.json}) and a SQLite database (\texttt{data/papers.db}),
enabling structured query, temporal tracking, and cross-paper analysis.

Outputs are produced in three formats: daily digests (markdown files
summarizing all papers collected in a given day), weekly synthesis reports
(aggregating and cross-referencing the week's collection), and topic-trend
matrices (tracking the proportion of papers assigned to each of the five
thematic axes over time). Dissemination to the research team is automated
via Telegram bot, with email as a secondary channel.

The pipeline runs daily via a scheduled cron job (08:00, weekdays),
ensuring that the review corpus is continuously refreshed as new
publications appear. This design choice — continuous collection rather
than a one-time snapshot — is methodologically significant: the health AI
field evolves too rapidly for a static search to capture the current state
of knowledge at the time of writing.

\subsection{Search Strategy and Information Sources}
\label{sec:search_strategy}

For the systematic review itself, papers were retrieved from the pipeline
database covering the period January 2020 to April 2026. The search was
replicated across four electronic databases to ensure comprehensive
coverage: \textbf{PubMed} (MEDLINE), \textbf{Semantic Scholar},
\textbf{arXiv}, and \textbf{Google Scholar}. The search query was
constructed around the five thematic axes described above, each comprising
controlled vocabulary terms and free-text keywords:

\begin{enumerate}[label=(\Roman*), itemsep=0.2em]
\item \textbf{AI Evaluation and Validation} — ``AI evaluation,'' ``AI
  validation,'' ``algorithmic auditing,'' ``model evaluation,'' ``bias
  detection,'' ``robustness testing,'' ``post-deployment monitoring,''
  ``continuous evaluation,'' ``explainability,'' ``safety evaluation.''

\item \textbf{Participatory Governance} — ``participatory governance,''
  ``stakeholder engagement,'' ``citizen participation,'' ``algorithmic
  accountability,'' ``democratic AI,'' ``community-centered AI.''

\item \textbf{Adaptive Regulation} — ``adaptive governance,'' ``agile
  regulation,'' ``responsive regulation,'' ``regulatory sandboxes,''
  ``living labs,'' ``evidence-based regulation,'' ``continuous compliance.''

\item \textbf{Evidence and Implementation} — ``evidence-based policy,''
  ``implementation science,'' ``knowledge translation,'' ``policy uptake,''
  ``evidence gaps.''

\item \textbf{Clinical Health AI} — ``clinical AI,'' ``medical AI,''
  ``clinical decision support,'' ``diagnostic AI,'' ``predictive modeling in
  healthcare,'' ``digital health,'' ``AI safety in healthcare.''

\end{enumerate}

Search strings were adapted for each database using Boolean operators
(AND/OR) and database-specific syntax. No language restrictions were applied
at the search stage; however, only English-language publications were
retained after abstract screening. Preprints and peer-reviewed articles
were both included given the rapid evolution of the field.

\subsection{Selection Criteria}
\label{sec:selection_criteria}

Studies were included if they met all of the following criteria:

\begin{itemize}[itemsep=0.2em]
\item \textbf{Population:} Human health AI systems deployed in clinical,
  sub-clinical, or community health settings; or AI systems evaluated for
  direct or indirect clinical application.
\item \textbf{Intervention:} Any AI-based intervention — including diagnostic
  AI, predictive modeling, clinical decision support systems, natural language
  processing for clinical text, wearable sensor analytics, and AI-assisted
  therapeutic applications.
\item \textbf{Outcomes:} Studies reporting on evaluation metrics, governance
  structures, implementation barriers, regulatory frameworks, ethical
  considerations, or societal impact of health AI.
\item \textbf{Study design:} Peer-reviewed empirical studies (RCTs, cohort
  studies, case studies, cross-sectional studies), systematic reviews,
  scoping reviews, conceptual papers, and policy analyses. Animal studies,
  purely computational studies without clinical application, and publications
  in predatory journals were excluded.
\item \textbf{Date range:} Published between 2020 and 2026.
\end{itemize}

\subsection{Quality Assessment}
\label{sec:quality_assessment}

Given the heterogeneous study designs included in this review, a modular
quality appraisal approach was adopted. For empirical clinical studies,
the \textbf{JBI Critical Appraisal Checklist} appropriate to the study
design was applied. For conceptual and policy papers, a bespoke quality
rubric assessed: (1) explicit statement of the research question or
analytical aim; (2) transparency in methods (search strategy, selection
criteria, data sources); (3) balanced representation of evidence including
acknowledgment of limitations; and (4) clarity and logical coherence of
argumentation. Studies were not excluded on the basis of quality alone;
rather, quality appraisal scores weighted the interpretive confidence
assigned to individual findings during synthesis.

\subsection{Data Extraction and Analysis}
\label{sec:data_extraction}

Data were extracted using a standardized extraction form capturing: study
metadata (authors, year, journal, country); study design; AI application
domain; evaluation methods and metrics; governance and regulatory mechanisms
described; implementation outcomes reported; and thematic codes assigned.
Extraction was performed by the primary reviewer and verified by a secondary
reviewer; discrepancies were resolved by discussion.

Analysis followed a \textbf{thematic synthesis} approach, combining inductive
coding of primary findings with a deductive framework organized around the
three analytical dimensions defined in the research question: (i) evaluation
methods and metrics; (ii) governance and regulatory frameworks; and
(iii) implementation outcomes and barriers. Results are reported narratively
organized by theme, with summary tables providing an overview of the
included studies and their key characteristics.

A PRISMA flow diagram is provided in Appendix~\ref{sec:prisma}. The review
protocol was designed and reported following a modified PRISMA framework
adapted for thematic synthesis of heterogeneous evidence.